\documentclass[12pt]{article}
\usepackage[utf8]{inputenc}
\usepackage[T1]{fontenc}
\usepackage{hyperref}
\title{Mixtepec Zapotec Language Resources for K-12 Curriculum Development}
\author{Compiled by Claude Sonnet 5.5 (claude-sonnet-5-5) for Andr\'as Kornai}
\date{}
\begin{document}
\maketitle

\section{Introduction}
Mixtepec Zapotec (ISO 639-3 zpm, Glottocode mixt1426) is an Otomanguean language of Oaxaca, Mexico, with about 7,000 speakers. Search found only a few academic descriptions and no school materials. No Talking Dictionary exists for this variety in the results; the Swarthmore Talking Dictionary hits are for Miahuatec Zapotec, and the ``Mixtepec-Mixtec Lexicon'' hit is for a different language.

\section{Primary Educational Resources}
No primers, textbooks or school programmes for Mixtepec Zapotec were found in search results.

\section{Linguistic and Reference Materials}
\begin{itemize}
\item Roger Reeck, \emph{A trilingual dictionary in (Mixtepec) Zapotec, English and Spanish}, 1991 MA thesis, Universidad de las Am\'ericas-Puebla. Bibliographic record only; no digital copy was found. \url{https://glottolog.org/resource/reference/id/40218}
\item Hunn, Baron, Reeck, Hern\'andez P\'erez and Silva Cruz, ``A Sketch of Mixtepec Zapotec Grammar'' (34 pages: phonology, pronouns, verbs, nouns). The only copy seen was on a file-download site of unclear provenance, so verify it against an academic source. \url{https://n.gst.elementfx.com/dl/a-sketch-of-mixtepec-zapotec-grammar}
\item Endangered Languages Project bibliography for Mixtepec Zapotec. \url{https://www.endangeredlanguages.com/lang/3124/bibliography?hl=zh-cn}
\end{itemize}

\section{Religious and Community Resources}
None specific to Mixtepec Zapotec were found in these searches.

\section{Digital and Audio Resources}
No audio or digital resources for Mixtepec Zapotec were found. The Miahuatec Zapotec (San Bartolom\'e Loxicha) Talking Dictionary at Swarthmore is a different variety and should not be mistaken for it.

\section{Curriculum Development Context}
The language is not known to be taught in schools, and available documentation is limited to a trilingual dictionary thesis and a short grammar sketch. Curriculum work would rely on those, on the Hunn et al. authors' fieldwork, and on new materials made with community speakers.

\end{document}